\documentclass[12pt, a4paper]{article}
\usepackage{../../../myconfig}

\begin{document}

\titlebox{Chủ đề: Oxyz}{Vecto và các phép toán vecto \\ trong không gian}{Oxyz}


% ==========================================
% I. LÝ THUYẾT TRỌNG TÂM: NHÓM 1 (TRANG 1: 4 Ý ĐẦU HÌNH TO, GIÃN THOÁNG ĐỀU)
% ==========================================
\noindent\textbf{\textcolor{deepblue}{\large I. TÓM TẮT TRỌNG TÂM (NHÓM 1: VECTƠ VÀ CÁC QUY TẮC)}}
\par\vspace{0.4cm}

\setlength{\abovedisplayskip}{7pt}
\setlength{\belowdisplayskip}{7pt}

% ------------------------------------------
% MỤC 1: VECTƠ TRONG KHÔNG GIAN
% ------------------------------------------
\noindent
\begin{minipage}[c]{0.63\linewidth}
    \textbf{\textcolor{amberorange}{1. Vectơ trong không gian và Vectơ-không}}
    \begin{itemize}\setlength{\itemsep}{4pt}\setlength{\parskip}{0pt}
        \item Là một đoạn thẳng có hướng: $\vec{AB}$, độ dài $|\vec{AB}| = AB$.
        \item $\vec{0}$ có điểm đầu trùng điểm cuối $(\vec{AA} = \vec{0})$, \newline
              có độ dài bằng $0$, cùng phương và cùng hướng với mọi vectơ.
    \end{itemize}
\end{minipage}\hfill
\begin{minipage}[c]{0.35\linewidth}
    \centering
    \begin{tikzpicture}[scale=1.0, font=\footnotesize]
        \coordinate (A) at (0,0);
        \coordinate (B) at (2.8,0);
        \coordinate (D) at (1.1,0.9);
        \coordinate (C) at ($(B)+(D)-(A)$);
        \coordinate (A1) at (0,1.5);
        \coordinate (B1) at ($(B)+(A1)-(A)$);
        \coordinate (C1) at ($(C)+(A1)-(A)$);
        \coordinate (D1) at ($(D)+(A1)-(A)$);

        \draw[dashed, softgray!80!black] (A)--(B)--(C)--(D)--cycle;
        \draw[dashed, softgray!80!black] (A1)--(B1)--(C1)--(D1)--cycle;
        \draw[dashed, softgray!80!black] (A)--(A1) (B)--(B1) (C)--(C1) (D)--(D1);

        \draw[->, >=stealth, line width=1.5pt, amberorange] (A) -- (C1) node[midway, above left] {$\vec{AB}$};
        \fill[deepblue] (A) circle (1.5pt) node[below left] {$A$};
        \fill[deepblue] (C1) circle (1.5pt) node[above right] {$B$};
    \end{tikzpicture}
\end{minipage}

\par\vspace{0.45cm}
\hrule height 0.4pt \relax
\par\vspace{0.45cm}

% ------------------------------------------
% MỤC 2: CÙNG PHƯƠNG, CÙNG HƯỚNG, BẰNG NHAU
% ------------------------------------------
\noindent
\begin{minipage}[c]{0.63\linewidth}
    \textbf{\textcolor{amberorange}{2. Cùng phương, cùng hướng, bằng nhau}}
    \begin{itemize}\setlength{\itemsep}{4pt}\setlength{\parskip}{0pt}
        \item \textbf{Cùng phương:} Giá song song hoặc trùng nhau.
        \item \textbf{Cùng hướng / Ngược hướng:} \newline
              Khi hai vectơ đã cùng phương.
        \item \textbf{Bằng nhau:} $\vec{a} = \vec{b} \iff$ \newline
              $\vec{a}, \vec{b} \text{ cùng hướng và } |\vec{a}| = |\vec{b}|$.
    \end{itemize}
\end{minipage}\hfill
\begin{minipage}[c]{0.35\linewidth}
    \centering
    \begin{tikzpicture}[scale=0.95, font=\scriptsize]
        \draw[->, >=stealth, deepblue, line width=1.1pt] (0, 1.2) -- (2.7, 1.2);
        \draw[->, >=stealth, deepblue, line width=1.1pt] (0, 0.75) -- (2.7, 0.75);
        \node[deepblue, font=\footnotesize, right] at (2.75, 0.98) {cùng hướng};

        \draw[->, >=stealth, accentred, line width=1.1pt] (0, 0.05) -- (2.7, 0.05);
        \draw[<-, >=stealth, accentred, line width=1.1pt] (0, -0.4) -- (2.7, -0.4);
        \node[accentred, font=\footnotesize, right] at (2.75, -0.18) {ngược hướng};

        \draw[->, >=stealth, amberorange, line width=1.5pt] (0, -1.1) -- (2.4, -1.1) node[midway, above] {$\vec{a}$};
        \draw[->, >=stealth, amberorange, line width=1.5pt] (0, -1.65) -- (2.4, -1.65) node[midway, above] {$\vec{b}$};
        \node[amberorange, font=\bfseries\footnotesize, right] at (2.55, -1.38) {$\vec{a} = \vec{b}$};
    \end{tikzpicture}
\end{minipage}

\par\vspace{0.45cm}
\hrule height 0.4pt \relax
\par\vspace{0.45cm}

% ------------------------------------------
% MỤC 3: QUY TẮC BA ĐIỂM (NỐI ĐUÔI)
% ------------------------------------------
\noindent
\begin{minipage}[c]{0.60\linewidth}
    \textbf{\textcolor{amberorange}{3. Quy tắc ba điểm (Nối đuôi hai vectơ)}}
    \begin{itemize}\setlength{\itemsep}{4pt}\setlength{\parskip}{0pt}
        \item Với ba điểm $A, B, C$ bất kì trong không gian:
        \[ \vec{AB} + \vec{BC} = \vec{AC} \]
        \item Điểm cuối của vectơ này là điểm đầu của vectơ kia.
    \end{itemize}
\end{minipage}\hfill
\begin{minipage}[c]{0.38\linewidth}
    \centering
    \begin{tikzpicture}[scale=1.1, font=\footnotesize]
        \coordinate (A) at (0, 0);
        \coordinate (B) at (1.8, 1.7);
        \coordinate (C) at (3.5, 0.2);

        \draw[->, >=stealth, deepblue, line width=1.2pt] (A) -- (B) node[midway, above left] {$\vec{AB}$};
        \draw[->, >=stealth, deepblue, line width=1.2pt] (B) -- (C) node[midway, above right] {$\vec{BC}$};
        \draw[->, >=stealth, amberorange, line width=1.5pt] (A) -- (C) node[midway, below] {$\vec{AC}$};

        \node[below left] at (A) {$A$};
        \node[above] at (B) {$B$};
        \node[right] at (C) {$C$};
    \end{tikzpicture}
\end{minipage}

\par\vspace{0.45cm}
\hrule height 0.4pt \relax
\par\vspace{0.45cm}

% ------------------------------------------
% MỤC 4: QUY TẮC HÌNH BÌNH HÀNH (CHUNG GỐC)
% ------------------------------------------
\noindent
\begin{minipage}[c]{0.60\linewidth}
    \textbf{\textcolor{amberorange}{4. Quy tắc hình bình hành (Tổng hai cạnh = Đường chéo)}}
    \begin{itemize}\setlength{\itemsep}{4pt}\setlength{\parskip}{0pt}
        \item Nếu $ABCD$ là hình bình hành:
        \[ \vec{AB} + \vec{AD} = \vec{AC} \]
        \item Áp dụng khi hai vectơ thành phần có cùng điểm xuất phát.
    \end{itemize}
\end{minipage}\hfill
\begin{minipage}[c]{0.38\linewidth}
    \centering
    \begin{tikzpicture}[scale=1.05, font=\footnotesize]
        \coordinate (A) at (0,0);
        \coordinate (B) at (3.0,0);
        \coordinate (D) at (1.3,1.4);
        \coordinate (C) at ($(B)+(D)-(A)$);

        \draw[dashed, softgray!80!black] (D) -- (C) -- (B);
        \draw[->, >=stealth, deepblue, line width=1.2pt] (A) -- (B) node[midway, below] {$\vec{AB}$};
        \draw[->, >=stealth, deepblue, line width=1.2pt] (A) -- (D) node[midway, above left] {$\vec{AD}$};
        \draw[->, >=stealth, amberorange, line width=1.5pt] (A) -- (C) node[midway, above] {$\vec{AC}$};

        \node[below left] at (A) {$A$};
        \node[below right] at (B) {$B$};
        \node[right] at (C) {$C$};
        \node[left] at (D) {$D$};
    \end{tikzpicture}
\end{minipage}

\newpage
% ==========================================
% TRANG 2: MỤC 5 & 6 (GIÃN THOÁNG ĐỀU TRANG, HÌNH TO ĐẸP)
% ==========================================
\vspace*{0.2cm}

% ------------------------------------------
% MỤC 5: QUY TẮC HÌNH HỘP
% ------------------------------------------
\noindent
\begin{minipage}[c]{0.58\linewidth}
    \textbf{\large\textcolor{amberorange}{5. Quy tắc hình hộp (Ba cạnh = Đường chéo)}}
    \begin{itemize}\setlength{\itemsep}{6pt}\setlength{\parskip}{0pt}
        \item Cho hình hộp $ABCD.A'B'C'D'$:
        \[ \vec{AB} + \vec{AD} + \vec{AA'} = \vec{AC'} \]
        \item Tổng ba vectơ xuất phát từ cùng một đỉnh bằng vectơ đường chéo xuất phát từ đỉnh đó.
        \item Tương tự tại các đỉnh khác:
        \[ \vec{BA} + \vec{BC} + \vec{BB'} = \vec{BD'} \]
        \[ \vec{C'A'} + \vec{C'C} + \vec{C'D'} = \vec{C'B} \]
    \end{itemize}
\end{minipage}\hfill
\begin{minipage}[c]{0.40\linewidth}
    \centering
    \begin{tikzpicture}[scale=1.1, font=\footnotesize]
        \coordinate (A)   at (0,0);
        \coordinate (B)   at (2.6,-0.5);
        \coordinate (D)   at (2.5,0.5);
        \coordinate (C)   at ($(B)+(D)-(A)$);
        
        \coordinate (A1)  at (0,2.2);
        \coordinate (B1)  at ($(B)+(A1)-(A)$);
        \coordinate (D1)  at ($(D)+(A1)-(A)$);
        \coordinate (C1)  at ($(C)+(A1)-(A)$);

        \draw[dashed, softgray!80!black] (A) -- (B) -- (C) -- (D) -- cycle;
        \draw[dashed, softgray!80!black] (A1) -- (B1) -- (C1) -- (D1) -- cycle;
        \draw[dashed, softgray!80!black] (A) -- (A1) (B) -- (B1) (C) -- (C1) (D) -- (D1);

        \draw[->, >=stealth, deepblue, line width=1.1pt] (A) -- (B) node[below left] {$\vec{AB}$};
        \draw[->, >=stealth, deepblue, line width=1.1pt] (A) -- (D) node[below right] {$\vec{AD}$};
        \draw[->, >=stealth, deepblue, line width=1.1pt] (A) -- (A1) node[left] {$\vec{AA'}$};
        \draw[->, >=stealth, amberorange, line width=1.6pt] (A) -- (C1) node[above right] {$\vec{AC'}$};
        \fill[deepblue] (A) circle (1.5pt) node[left] {$A$};
    \end{tikzpicture}
\end{minipage}

\par\vspace{1.0cm}
\hrule height 0.5pt \relax
\par\vspace{1.0cm}

% ------------------------------------------
% MỤC 6: NHÂN VECTƠ VỚI MỘT SỐ
% ------------------------------------------
\noindent
\begin{minipage}[c]{0.63\linewidth}
    \textbf{\large\textcolor{amberorange}{6. Nhân vectơ với một số (Điều kiện cùng phương)}}
    \begin{itemize}\setlength{\itemsep}{5pt}\setlength{\parskip}{0pt}
        \item \textbf{Độ dài:} $|k\vec{a}| = |k| \cdot |\vec{a}|$.
        \item \textbf{Hướng:}
        \begin{itemize}\setlength{\itemsep}{2pt}
            \item Cùng hướng với $\vec{a}$ khi $k > 0$.
            \item Ngược hướng với $\vec{a}$ khi $k < 0$.
        \end{itemize}
        \item \textbf{Điều kiện cùng phương:}
        \[ \vec{a} \parallel \vec{b} \iff \vec{a} = k\vec{b} \quad (k \in \mathbb{R},\ \vec{b} \neq \vec{0}) \]
        \item Ba điểm phân biệt $A, B, C$ thẳng hàng khi và chỉ khi tồn tại số $k$ sao cho $\vec{AB} = k\vec{AC}$.
    \end{itemize}
\end{minipage}\hfill
\begin{minipage}[c]{0.35\linewidth}
    \centering
    \begin{tikzpicture}[scale=1.0, font=\scriptsize]
        % Đường gióng mốc gốc
        \draw[dashed, softgray!80!black, line width=0.8pt] (0, 1.4) -- (0, -1.3);
        \node[softgray!80!black, font=\tiny, above] at (0, 1.4) {Gốc};

        % 1. Vectơ a (độ dài 1.1cm)
        \draw[->, >=stealth, deepblue, line width=1.3pt] (0, 0.9) -- (1.1, 0.9);
        \node[deepblue, font=\bfseries\scriptsize, right] at (1.15, 0.9) {$\vec{a}$};

        % 2. Vectơ 2a: k = 2 > 0 (độ dài 2.2cm)
        \draw[->, >=stealth, deepblue, line width=1.3pt] (0, 0.0) -- (2.2, 0.0);
        \node[deepblue, font=\bfseries\scriptsize, right] at (2.25, 0.0) {$2\vec{a}\ (k>0)$};

        % 3. Vectơ -a: k = -1 < 0 (sang trái 1.1cm)
        \draw[->, >=stealth, accentred, line width=1.3pt] (0, -0.9) -- (-1.1, -0.9);
        \node[accentred, font=\bfseries\scriptsize, left] at (-1.15, -0.9) {$-\vec{a}\ (k<0)$};
        
        % Chấm tròn gốc
        \fill[deepblue] (0, 0.9) circle (1.6pt);
        \fill[deepblue] (0, 0.0) circle (1.6pt);
        \fill[accentred] (0, -0.9) circle (1.6pt);
    \end{tikzpicture}
\end{minipage}

\vfill
\newpage
% ==========================================
% II. VÍ DỤ MINH HỌA TRÊN LỚP (NHÓM 1)
% ==========================================
\noindent\textbf{\textcolor{deepblue}{\large II. VÍ DỤ MINH HỌA TRÊN LỚP (NHÓM 1)}}
\par\vspace{0.25cm}

% Ví dụ 1: Ý 1 & Ý 2
\questionLinesManual{0}{1}{Cho tứ diện đều $ABCD$ có cạnh bằng $a$.
\par
\textbf{a)} Tính độ dài vectơ $\vec{u} = \vec{AB} - \vec{AC}$.
\par
\textbf{b)} Gọi $M, N$ lần lượt là trung điểm của $AB, AC$. Tìm số thực $k$ sao cho $\vec{MN} = k\vec{BC}$ và tính độ dài $|\vec{MN}|$.}{}
\drawdashlinesfull

% Ví dụ 2: Ý 2 & Ý 6
\questionLinesManual{0}{2}{Cho hình bình hành $ABCD$ có tâm $O$. 
\par
\textbf{a)} Tìm tất cả các vectơ có điểm đầu và điểm cuối là các đỉnh của hình bình hành thỏa mãn bằng vectơ $\vec{AB}$.
\par
\textbf{b)} Tìm số thực $k$ thỏa mãn đẳng thức $\vec{AO} = k\vec{CA}$.}{}
\drawdashlinesfull

% Ví dụ 3: Ý 3 & Ý 4
\questionLinesManual{0}{3}{Cho tứ diện $ABCD$.
\par
\textbf{a)} Rút gọn biểu thức tổng: $\vec{x} = \vec{AB} + \vec{BC} + \vec{CD} + \vec{DA}$.
\par
\textbf{b)} Rút gọn biểu thức: $\vec{y} = \vec{AC} - \vec{AD} + \vec{BD}$ về dạng một vectơ có điểm đầu là $B$.}{}
\drawdashlinesfull

% Ví dụ 4: Ý 4 & Ý 6
\questionLinesManual{0}{4}{Cho hình chóp $S.ABCD$ có đáy $ABCD$ là hình bình hành tâm $O$.
\par
\textbf{a)} Rút gọn biểu thức vectơ: $\vec{u} = \vec{SA} + \vec{SC} - (\vec{SB} + \vec{SD})$.
\par
\textbf{b)} Cho biết $\vec{SA} + \vec{SC} = k\vec{SO}$. Hãy xác định giá trị của số thực $k$.}{}
\drawdashlinesfull

% Ví dụ 5: Ý 5
\questionLinesManual{0}{5}{Cho hình hộp chữ nhật $ABCD.A'B'C'D'$ có các kích thước $AB = 3$, $AD = 4$, $AA' = 5$.
\par
\textbf{a)} Rút gọn tổng ba vectơ $\vec{v} = \vec{AB} + \vec{AD} + \vec{AA'}$ thành một vectơ đường chéo duy nhất của hình hộp.
\par
\textbf{b)} Tính độ dài của vectơ $\vec{v}$.}{}
\drawdashlinesfull

% Ví dụ 6: Ý 5 & Ý 6
\questionLinesManual{0}{6}{Cho hình hộp $ABCD.A'B'C'D'$. Gọi $M$ là điểm xác định bởi đẳng thức $\vec{AM} = \dfrac{1}{3}\vec{AC'}$. 
\par
Biểu diễn vectơ $\vec{AM}$ theo ba vectơ $\vec{a} = \vec{AB}$, $\vec{b} = \vec{AD}$, $\vec{c} = \vec{AA'}$ dưới dạng $\vec{AM} = x\vec{a} + y\vec{b} + z\vec{c}$. Tính tổng $T = x + y + z$.}{}
\drawdashlinesfull

\newpage




% ==========================================
% I. LÝ THUYẾT TRỌNG TÂM: NHÓM 2 (GIÃN ĐỀU 1 TRANG, HÌNH TO RÕ)
% ==========================================
\noindent\textbf{\textcolor{deepblue}{\large I. TÓM TẮT TRỌNG TÂM (NHÓM 2: CÁC HỆ THỨC VECTƠ)}}
\par\vspace{0.35cm}

\setlength{\abovedisplayskip}{6pt}
\setlength{\belowdisplayskip}{6pt}

% ------------------------------------------
% MỤC 7: HỆ THỨC TRUNG ĐIỂM
% ------------------------------------------
\noindent
\begin{minipage}[c]{0.65\linewidth}
    \textbf{\textcolor{amberorange}{7. Hệ thức trung điểm}}
    \begin{itemize}\setlength{\itemsep}{2pt}\setlength{\parskip}{0pt}
        \item Điểm $I$ là trung điểm của đoạn thẳng $AB$:
        \[ \vec{IA} + \vec{IB} = \vec{0} \]
        \item Với điểm $M$ bất kì trong không gian:
        \[ \vec{MA} + \vec{MB} = 2\vec{MI} \]
    \end{itemize}
\end{minipage}\hfill
\begin{minipage}[c]{0.33\linewidth}
    \centering
    \begin{tikzpicture}[scale=0.85, font=\footnotesize]
        \coordinate (A) at (0,0);
        \coordinate (B) at (3.2,0);
        \coordinate (I) at (1.6,0);
        \coordinate (M) at (1.6,1.8);

        \draw[thick, deepblue] (M) -- (A) -- (B) -- cycle;
        \draw[->, >=stealth, amberorange, line width=1.4pt] (M) -- (I) node[midway, right] {$\vec{MI}$};

        \fill[deepblue] (A) circle (1.3pt) node[below left] {$A$};
        \fill[deepblue] (B) circle (1.3pt) node[below right] {$B$};
        \fill[deepblue] (I) circle (1.3pt) node[below] {$I$};
        \fill[deepblue] (M) circle (1.3pt) node[above] {$M$};
    \end{tikzpicture}
\end{minipage}

\par\vspace{0.25cm}
\hrule height 0.4pt \relax
\par\vspace{0.25cm}

% ------------------------------------------
% MỤC 8: TRỌNG TÂM TAM GIÁC
% ------------------------------------------
\noindent
\begin{minipage}[c]{0.65\linewidth}
    \textbf{\textcolor{amberorange}{8. Trọng tâm tam giác}}
    \begin{itemize}\setlength{\itemsep}{2pt}\setlength{\parskip}{0pt}
        \item Điểm $G$ là trọng tâm của tam giác $ABC$:
        \[ \vec{GA} + \vec{GB} + \vec{GC} = \vec{0} \]
        \item Với điểm $M$ bất kì trong không gian:
        \[ \vec{MA} + \vec{MB} + \vec{MC} = 3\vec{MG} \]
    \end{itemize}
\end{minipage}\hfill
\begin{minipage}[c]{0.33\linewidth}
    \centering
    \begin{tikzpicture}[scale=0.85, font=\footnotesize]
        \coordinate (A) at (1.5, 2.0);
        \coordinate (B) at (0, 0);
        \coordinate (C) at (3.0, 0.2);
        \coordinate (G) at ($1/3*(A) + 1/3*(B) + 1/3*(C)$);

        \draw[dashed, softgray!80!black] (A) -- (B) -- (C) -- cycle;
        \draw[->, >=stealth, amberorange, line width=1.3pt] (G) -- (A);
        \draw[->, >=stealth, amberorange, line width=1.3pt] (G) -- (B);
        \draw[->, >=stealth, amberorange, line width=1.3pt] (G) -- (C);

        \fill[deepblue] (A) circle (1.3pt) node[above] {$A$};
        \fill[deepblue] (B) circle (1.3pt) node[below left] {$B$};
        \fill[deepblue] (C) circle (1.3pt) node[below right] {$C$};
        \fill[amberorange] (G) circle (1.5pt) node[right] {$G$};
    \end{tikzpicture}
\end{minipage}

\par\vspace{0.25cm}
\hrule height 0.4pt \relax
\par\vspace{0.25cm}

% ------------------------------------------
% MỤC 9: TRỌNG TÂM TỨ DIỆN
% ------------------------------------------
\noindent
\begin{minipage}[c]{0.65\linewidth}
    \textbf{\textcolor{amberorange}{9. Trọng tâm tứ diện}}
    \begin{itemize}\setlength{\itemsep}{2pt}\setlength{\parskip}{0pt}
        \item Điểm $G$ là trọng tâm của tứ diện $ABCD$:
        \[ \vec{GA} + \vec{GB} + \vec{GC} + \vec{GD} = \vec{0} \]
        \item Với điểm $M$ bất kì trong không gian:
        \[ \vec{MA} + \vec{MB} + \vec{MC} + \vec{MD} = 4\vec{MG} \]
    \end{itemize}
\end{minipage}\hfill
\begin{minipage}[c]{0.33\linewidth}
    \centering
    \begin{tikzpicture}[scale=0.82, font=\footnotesize]
        \coordinate (A) at (1.5, 2.2);
        \coordinate (B) at (0, 0);
        \coordinate (D) at (1.0, 0.7);
        \coordinate (C) at (3.0, 0.2);
        \coordinate (G) at ($0.25*(A) + 0.25*(B) + 0.25*(C) + 0.25*(D)$);

        \draw[dashed, softgray!80!black] (B) -- (D) -- (C) (A) -- (D);
        \draw[thick, softgray!80!black] (A) -- (B) -- (C) -- cycle;

        \draw[->, >=stealth, amberorange, line width=1.2pt] (G) -- (A);
        \draw[->, >=stealth, amberorange, line width=1.2pt] (G) -- (B);
        \draw[->, >=stealth, amberorange, line width=1.2pt] (G) -- (C);
        \draw[->, >=stealth, amberorange, line width=1.2pt] (G) -- (D);

        \fill[deepblue] (A) circle (1.3pt) node[above] {$A$};
        \fill[deepblue] (B) circle (1.3pt) node[below left] {$B$};
        \fill[deepblue] (D) circle (1.3pt) node[left] {$D$};
        \fill[deepblue] (C) circle (1.3pt) node[right] {$C$};
        \fill[amberorange] (G) circle (1.5pt) node[above right] {$G$};
    \end{tikzpicture}
\end{minipage}

\par\vspace{0.25cm}
\hrule height 0.4pt \relax
\par\vspace{0.25cm}

% ------------------------------------------
% MỤC 10: TÍCH VÔ HƯỚNG VÀ GÓC GIỮA HAI VECTƠ
% ------------------------------------------
\noindent
\begin{minipage}[c]{0.65\linewidth}
    \textbf{\textcolor{amberorange}{10. Tích vô hướng của hai vectơ}}
    \begin{itemize}\setlength{\itemsep}{2pt}\setlength{\parskip}{0pt}
        \item Góc giữa hai vectơ $\vec{a}$ và $\vec{b}$ thỏa mãn: $0^\circ \le (\vec{a}, \vec{b}) \le 180^\circ$.
        \item Công thức tích vô hướng:
        \[ \vec{a} \cdot \vec{b} = |\vec{a}| \cdot |\vec{b}| \cdot \cos(\vec{a}, \vec{b}) \]
    \end{itemize}
\end{minipage}\hfill
\begin{minipage}[c]{0.33\linewidth}
    \centering
    \begin{tikzpicture}[scale=1.0, font=\small]
    \coordinate (O) at (0,0);
    \coordinate (B) at (2.2,0);
    \coordinate (A) at (1.5,1.5); % góc 45 độ

    % Vẽ hai vectơ a và b
    \draw[->, >=stealth, deepblue, line width=1.2pt] (O) -- (B) node[below] {$\vec{b}$};
    \draw[->, >=stealth, deepblue, line width=1.2pt] (O) -- (A) node[above] {$\vec{a}$};

    % Vẽ cung tròn góc phi
    \draw[amberorange, line width=1.2pt] (0.6,0) arc (0:45:0.6);
    \node[amberorange, font=\footnotesize] at (0.85, 0.3) {$\varphi$};

    % Điều kiện góc được đẩy sang phải rõ ràng, không bị đè nét
    \node[deepblue, font=\scriptsize, anchor=west] at (1.1, 0.85) {$0^\circ \le \varphi \le 180^\circ$};
\end{tikzpicture}
\end{minipage}

\par\vspace{0.25cm}
\hrule height 0.4pt \relax
\par\vspace{0.25cm}

% ------------------------------------------
% MỤC 11: VUÔNG GÓC VÀ BÌNH PHƯƠNG VÔ HƯỚNG
% ------------------------------------------
\noindent
\begin{minipage}[c]{0.65\linewidth}
    \textbf{\textcolor{amberorange}{11. Vuông góc và bình phương vô hướng}}
    \begin{itemize}\setlength{\itemsep}{2pt}\setlength{\parskip}{0pt}
        \item Hai vectơ vuông góc khi và chỉ khi tích vô hướng bằng $0$:
        \[ \vec{a} \perp \vec{b} \iff \vec{a} \cdot \vec{b} = 0 \]
        \item Bình phương vô hướng bằng bình phương độ dài:
        \[ \vec{a}^2 = |\vec{a}|^2 \]
    \end{itemize}
\end{minipage}\hfill
\begin{minipage}[c]{0.33\linewidth}
    \centering
    \begin{tikzpicture}[scale=0.85, font=\footnotesize]
        \coordinate (O) at (0,0);
        \draw[->, >=stealth, deepblue, thick] (O) -- (2.5,0) node[below] {$\vec{b}$};
        \draw[->, >=stealth, deepblue, thick] (O) -- (0,2.0) node[left] {$\vec{a}$};
        \draw[amberorange, line width=1.1pt] (0.4,0) -- (0.4,0.4) -- (0,0.4);
        \node[amberorange, font=\small\bfseries] at (1.4, 1.1) {$\vec{a}\cdot\vec{b} = 0$};
    \end{tikzpicture}
\end{minipage}

\newpage




% ==========================================
% II. VÍ DỤ MINH HỌA TRÊN LỚP (NHÓM 2)
% ==========================================
\noindent\textbf{\textcolor{deepblue}{\large II. VÍ DỤ MINH HỌA TRÊN LỚP (NHÓM 2)}}
\par\vspace{0.25cm}

% Ví dụ 1: Ý 7 (Hệ thức trung điểm - Tính toán cụ thể)
\questionLinesManual{0}{1}{Cho tứ diện $ABCD$. Gọi $M, N$ lần lượt là trung điểm của các cạnh $AB$ và $CD$.
\par
\textbf{a)} Rút gọn biểu thức tổng: $\vec{u} = \vec{MA} + \vec{MB}$.
\par
\textbf{b)} Biết rằng $\vec{AC} + \vec{BD} = k\vec{MN}$. Hãy tìm giá trị của số thực $k$.}{}
\drawdashlinesfull

% Ví dụ 2: Ý 8 (Trọng tâm tam giác - Phân tích hệ số)
\questionLinesManual{0}{2}{Cho hình chóp $S.ABC$. Gọi $G$ là trọng tâm của tam giác $ABC$.
\par
\textbf{a)} Phân tích vectơ $\vec{SG}$ theo ba vectơ $\vec{SA}$, $\vec{SB}$, $\vec{SC}$ dưới dạng $\vec{SG} = x\vec{SA} + y\vec{SB} + z\vec{SC}$.
\par
\textbf{b)} Tính giá trị biểu thức $P = x + y + z$.}{}
\drawdashlinesfull

% Ví dụ 3: Ý 9 (Trọng tâm tứ diện - Tính tổng hệ số)
\questionLinesManual{0}{3}{Cho tứ diện $ABCD$ có trọng tâm $G$. Lấy điểm $O$ bất kì trong không gian thỏa mãn:
$$\vec{OG} = m(\vec{OA} + \vec{OB} + \vec{OC} + \vec{OD})$$
\par
\textbf{a)} Tìm giá trị của số thực $m$.
\par
\textbf{b)} Khi điểm $O$ trùng với đỉnh $A$, biểu diễn vectơ $\vec{AG}$ theo ba vectơ $\vec{AB}$, $\vec{AC}$ và $\vec{AD}$.}{}
\drawdashlinesfull

% Ví dụ 4: Ý 10 (Tích vô hướng & Tính góc cụ thể)
\questionLinesManual{0}{4}{Cho hình chóp tam giác đều $S.ABC$ có tất cả các cạnh đều bằng $a$.
\par
\textbf{a)} Xác định góc giữa hai vectơ $\vec{AB}$ và $\vec{AC}$, từ đó tính tích vô hướng $\vec{AB} \cdot \vec{AC}$ theo $a$.
\par
\textbf{b)} Tính giá trị tích vô hướng $\vec{SA} \cdot \vec{BC}$.}{}
\drawdashlinesfull

% Ví dụ 5: Ý 11 (Bình phương vô hướng & Tính độ dài)
\questionLinesManual{0}{5}{Cho hai vectơ $\vec{a}$ và $\vec{b}$ thỏa mãn $|\vec{a}| = 2$, $|\vec{b}| = 3$ và $(\vec{a}, \vec{b}) = 60^\circ$.
\par
\textbf{a)} Tính tích vô hướng $\vec{a} \cdot \vec{b}$.
\par
\textbf{b)} Tính độ dài của vectơ tổng $\vec{u} = 2\vec{a} - \vec{b}$.}{}
\drawdashlinesfull

\end{document}